\chapter{术语与缩写}
\begin{longtable}{p{.20\linewidth}p{.70\linewidth}}
\toprule 名称 & 本报告含义\\\midrule\endhead
ROS2 & 机器人操作系统第二代，提供节点通信、消息和时间等接口。\\
Bridge & 通信桥，连接仿真中的关节状态与外部ROS2命令。\\
QoS / DDS & 服务质量配置与数据分发通信机制；会影响可靠性、队列与传输。\\
PhysX / MuJoCo & 用于刚体动力学、碰撞与接触求解的物理引擎。\\
IMU & 惯性测量单元，常输出三轴加速度、三轴角速度等信息。\\
EMA & 指数移动平均，只使用当前与历史数据的递推滤波。\\
Butterworth & 巴特沃斯滤波器，通带幅频响应平坦；本项目采用因果实现。\\
Allan偏差 & 观察不同平均时间下的信号波动，单独曲线不等于完成设备标定。\\
群延迟 & 信号经过滤波器后不同频率分量的时间延迟特性。\\
数据年龄 & 控制器读取时，所用测量距离其源时间戳过去的时间。\\
RMSE & 均方根误差；需要同时注明参考、窗口和物理单位。\\
PD / PID & 比例、微分反馈；PID另外加入积分反馈。\\
导纳控制 & 根据受力误差生成运动修正；本项目主要研究平面法向响应。\\
逆运动学 & 从末端位置姿态求关节参考；微分逆解使用雅可比计算连续增量。\\
雅可比 & 将关节速度映射为末端或任务空间速度的局部线性关系。\\
LIPM & 线性倒立摆模型，将机器人简化为恒高质点与地面支撑点。\\
CoM & 质心，本报告同时关注位置、速度及加速度目标。\\
CoP & 压力中心，由实际地面接触力和力矩计算。\\
ZMP & 零力矩点；在恒高、零质心角动量变化等近似下与CoP对应。\\
DCM & 发散分量，将质心位置和速度组合，用于分析不稳定运动分量。\\
LQR & 线性二次调节器，根据状态误差与输入代价求反馈增益。\\
QP / WBC & 二次规划／全身控制，协调任务目标与动力学、接触和执行器约束。\\
接触扳手 & 三维力与力矩的组合；须说明参考点、坐标和分量顺序。\\
浮动基座 & 基座在空间中自由运动，通常没有直接驱动，需要通过关节和接触稳定。\\
双支撑 & 两脚共同承重的阶段；计划分配与实际接触反力分别记录。\\
准静态 & 运动较慢，以逐步转移重心和支撑为主；不代表动态高速步态。\\
OSQP / Clarabel & 本项目使用的数值优化求解器，不是两套不同机器人模型。\\
SO(3) & 三维旋转的数学表示；用于姿态误差与空间任务。\\
P50 / P95 / P99 & 第50、95、99百分位，描述分布及尾部，需注明样本量。\\
UDP / RTT & 用户数据报协议／往返时间；同机轮询耗时不等同于纯网络延迟。\\
JSON / CSV / NPZ & 结构化文本、表格与压缩数组文件，用于配置和保存结果。\\
SHA256 & 文件内容摘要，用于核对源码、模型、素材和输出是否变化。\\
VLA / 动作块 & 视觉语言动作模型／多步动作输出；本报告接口实验使用解析轨迹，不运行VLA。\\
\bottomrule\end{longtable}
